% ============================================================================
% 1. INTRODUCTION — Journal version (expanded)
% ============================================================================
\section{Introduction}
\label{sec:introduction}

Text-to-SQL benchmarks like Spider~\cite{yu2018spider} and BIRD~\cite{li2023bird}
have driven substantial progress in LLM-based database interaction.
Yet evaluating whether an LLM can \emph{generate} a SQL query against a schema
is distinct from evaluating whether it \emph{understands} that schema's structure.
As newer benchmarks like Spider~2.0~\cite{lei2024spider2} demonstrate,
enterprise-scale SQL workflows with complex multi-step reasoning remain
largely unsolved (best models achieve $\sim$17--21\% on Spider~2.0):
but even Spider~2.0 targets query generation, not structural topology reasoning.
SQL generation represents only one facet of database understanding.

Data warehouse environments require a fundamentally different class of
reasoning: understanding the \emph{graph topology} of schema relationships.
A data engineer maintaining an enterprise warehouse must routinely answer
structural questions such as:

\begin{itemize}
  \item \emph{``If the \texttt{Person} table changes, which downstream
    tables are affected through data lineage?''}: impact analysis
    requiring forward lineage traversal composed with reverse FK propagation.
  \item \emph{``Which tables form disconnected silos in our schema graph?''}
   : connectivity analysis requiring weakly connected component detection.
  \item \emph{``What is the shortest FK path between two tables?''}:
    routing requiring Dijkstra or BFS over the FK subgraph.
  \item \emph{``How many downstream rows are affected if we delete
    records from a staging table?''}: value-level impact analysis
    requiring interactive access to row-level lineage data.
\end{itemize}

These tasks require \textbf{graph traversal}: BFS, connected component
detection, multi-hop path finding: not SQL generation.  Yet no existing
benchmark evaluates LLMs on these capabilities.  Existing graph reasoning
benchmarks (NLGraph~\cite{wang2023nlgraph}, GraphArena~\cite{grapharena2024},
GraCoRe~\cite{gracore2024}) use synthetic or generic graphs and do not
capture the domain-specific challenges of data warehouse schemas with
\emph{heterogeneous edge types} (foreign key vs.\ lineage) and
\emph{mixed-layer topologies} (OLTP source tables connected to DW
destination tables via derived-from edges).

\paragraph{Contributions.} We make the following contributions:

\begin{enumerate}
  \item \textbf{DW-Bench}: A two-tier benchmark of 1,046 schema-level and
    433 value-level automatically-generated questions across 5 data warehouse
    schemas (4 real-world, 1 synthetic), organized into 3 categories
    (lineage impact, schema routing, silo detection), 13 Tier-1 subtypes,
    8 Tier-2 subtypes, and 3 difficulty levels.

  \item \textbf{Six baselines} spanning the text-only to agentic spectrum
    (Flat Text, Vector-RAG, Graph-Augmented, Tool-Use, ReAct-Code, and
    Oracle), evaluated with three frontier models (Gemini~2.5 Flash~\cite{gemini2024},
    DeepSeek-V3~\cite{deepseek2024v3}, and Qwen2.5-72B~\cite{qwen2024}), providing a systematic analysis of how different
    context injection strategies affect graph topology reasoning.

  \item \textbf{A two-tier evaluation framework}: Tier~1 tests schema-level
    topology reasoning (FK paths, lineage chains, components), while Tier~2
    tests value-level provenance tracing and impact analysis that requires
    interactive data access beyond context window limits.

  \item \textbf{An obfuscation protocol} that randomizes table names,
    enabling controlled evaluation of memorization versus genuine topology
    understanding, revealing that tool-based approaches are nearly
    obfuscation-invariant while static approaches lose 9--32\% EM.

  \item \textbf{Syn-Logistics}: A synthetic dataset with $n\geq20$ per
    subtype, addressing statistical power limitations of the real-world
    corpus and enabling reliable per-subtype analysis.

  \item \textbf{Open-source release} of all datasets, evaluation code,
    baseline implementations, and per-question results to enable
    community-driven analysis and benchmarking.
\end{enumerate}

